\documentclass[12pt]{article}
\usepackage{geometry}
\geometry{margin=1in}

% ----------------------------------------------------------
%  Linux Penguin Theme — Based on Recommended Colors
% ----------------------------------------------------------

% --- COLOR SETUP ---
\usepackage{xcolor}

\definecolor{cream}{RGB}{246,242,219}
\definecolor{teamred}{RGB}{209,57,64}
\definecolor{navy}{RGB}{40,46,87}
\definecolor{softwhite}{RGB}{250,250,230}
\definecolor{accentyellow}{RGB}{240,200,60}

% Extra utility colors
\definecolor{muted}{RGB}{120,125,140}
\definecolor{darkcream}{RGB}{230,225,205}

% ----------------------------------------------------------
%  FONT SETUP (Optional but Recommended)
%  Works best with XeLaTeX or LuaLaTeX
% ----------------------------------------------------------
\usepackage{fontspec}

% Main document font (Google fonts — upload to Overleaf)
% If you don't upload fonts, comment these two lines.
\setmainfont{SourceSans3}[
  Path=./,
  ItalicFont  = *-Italic,
  BoldFont    = *-Bold
]

\newfontfamily\headingfont{ScienceGothic}[
  Path=./,
]

% Font Awesome for minimal icons
\usepackage{fontawesome5}

% Fix for missing tree characters
\usepackage{pmboxdraw}
\usepackage{textcomp}

% ----------------------------------------------------------
%  SECTION STYLE
% ----------------------------------------------------------
\usepackage{titlesec}

% Section
\titleformat{\section}
  {\Large\headingfont\bfseries\color{navy}}
  {\thesection}
  {1em}
  {\rule{\linewidth}{1pt}\vspace{0.5em}\\}

% Subsection
\titleformat{\subsection}
  {\large\headingfont\bfseries\color{teamred}}
  {\thesubsection}
  {0.8em}
  {}

% Subsubsection
\titleformat{\subsubsection}
  {\bfseries\color{navy}}
  {\thesubsubsection}
  {0.6em}
  {}

% ----------------------------------------------------------
%  PAGE STYLE
% ----------------------------------------------------------
\usepackage{fancyhdr}
\pagestyle{fancy}

\fancyhf{} % clear all

% Fix headheight warning
\setlength{\headheight}{15.5pt}

\lhead{\textcolor{navy}{\headingfont \courseTitle}}
\rhead{\textcolor{teamred}{\thepage}}
\renewcommand{\headrule}{\color{teamred}\hrule height 1pt}

% ----------------------------------------------------------
%  HYPERLINK STYLE
% ----------------------------------------------------------
\usepackage{hyperref}

\hypersetup{
  colorlinks = true,
  linkcolor  = teamred,
  urlcolor   = navy,
  citecolor  = accentyellow
}

% ----------------------------------------------------------
%  CUSTOM BOX ENVIRONMENTS
% ----------------------------------------------------------
\usepackage{tcolorbox}
\tcbuselibrary{skins,breakable}

% Info box
\newtcolorbox{infobox}{
  enhanced,
  breakable,
  colback   = softwhite,
  colframe  = navy,
  coltext   = navy,
  boxrule   = 1pt,
  arc       = 4pt,
  left=6pt, right=6pt, top=6pt, bottom=6pt
}

% Warning box
\newtcolorbox{warningbox}{
  enhanced,
  breakable,
  colback   = accentyellow!30,
  colframe  = accentyellow!80!navy,
  coltext   = navy,
  boxrule   = 1pt,
  arc       = 4pt,
  left=6pt, right=6pt, top=6pt, bottom=6pt
}

% Highlight box
\newtcolorbox{highlightbox}{
  enhanced,
  breakable,
  colback   = teamred!10,
  colframe  = teamred,
  coltext   = navy,
  boxrule   = 1pt,
  arc       = 4pt,
  left=6pt, right=6pt, top=6pt, bottom=6pt
}

% ----------------------------------------------------------
%  ITEMIZE / ENUMERATE STYLE
% ----------------------------------------------------------
\usepackage{enumitem}

\setlist[itemize]{
  label=\textcolor{teamred}{\large\textbullet},
  leftmargin=1.2em,
  itemsep=4pt
}

\setlist[enumerate]{
  label=\textcolor{navy}{\arabic*.},
  leftmargin=1.4em,
  itemsep=4pt
}

% ----------------------------------------------------------
%  TITLE STYLE
% ----------------------------------------------------------
\usepackage{titling}

\pretitle{
  \begin{center}
  \headingfont\bfseries\LARGE\color{navy}
}
\posttitle{
  \end{center}
}

\preauthor{
  \begin{center}
  \color{teamred}
}
\postauthor{
  \end{center}
}

% ----------------------------------------------------------
%  OPTIONAL — SOFT PAGE BACKGROUND
%  Comment out if you want white pages
% ----------------------------------------------------------
\usepackage{pagecolor}
\pagecolor{cream}
\color{navy}

% ----------------------------------------------------------
% CUSTOM COMMANDS FOR TEMPLATE
% ----------------------------------------------------------
\newcommand{\courseTitle}{Managing Disks, RAM Usage, and Swap on Linux}
\newcommand{\courseDate}{02-03-2026}
\newcommand{\authorName}{Khaled Mahfouz}
\newcommand{\contactInfo}{\texttt{@khaledmhfz}}

% Minimal icon commands - FIXED: Using only valid Font Awesome 5 icon names
\newcommand{\penguinicon}{\faLinux}
\newcommand{\diskon}{\faHdd}              % Valid
\newcommand{\ramicon}{\faMicrochip}        % Valid
\newcommand{\swapicon}{\faArrowsAltH}      % FIXED: Replaced \faExchangeAlt with \faArrowsAltH
\newcommand{\zramicon}{\faCompress}        % Valid
\newcommand{\warningicon}{\faExclamationTriangle}
\newcommand{\tipsicon}{\faLightbulb}
\newcommand{\arrowicon}{\faArrowRight}
\newcommand{\bookicon}{\faBook}
\newcommand{\videoicon}{\faVideo}
\newcommand{\labicon}{\faFlask}
\newcommand{\rocketicon}{\faRocket}
\newcommand{\checkicon}{\faCheck}
\newcommand{\mounticon}{\faFolderOpen}
\newcommand{\externalicon}{\faExternalLink} % FIXED: Replaced external-link with \faExternalLink

% Tree drawing characters
\newcommand{\treebranch}{├── }
\newcommand{\treeleaf}{└── }

% ----------------------------------------------------------
% END OF THEME
% ----------------------------------------------------------

\begin{document}

% Title with minimal space
\begin{center}
\headingfont\Huge\bfseries\color{navy}
\diskon\ \courseTitle\\[0.5em]
\Large\color{teamred}
\courseDate
\end{center}

\vspace{1em}
\begin{center}
\color{teamred}\rule{0.8\textwidth}{1pt}
\end{center}

\section*{1 \faQuestionCircle\ Why This Session Matters}

Before we touch commands, understand this:

\begin{itemize}
    \item Your \textbf{disk} stores data permanently.
    \item Your \textbf{RAM} is fast temporary memory.
    \item \textbf{Swap} is emergency backup memory.
    \item \textbf{zram} is compressed RAM swap (smart trick for low-memory systems).
\end{itemize}

If you understand this session, you'll:

\begin{itemize}
    \item Mount extra drives safely
    \item Fix broken boot issues
    \item Add a second disk like a pro
    \item Understand RAM pressure
    \item Boost low-RAM machines using zram
\end{itemize}

\begin{infobox}
This is not theory. This is real Linux power.
\end{infobox}

\section*{2 \faSearch\ Checking If \texttt{zram-tools} Is Installed}

On Debian/Ubuntu/Mint, zram is handled by:

\begin{highlightbox}
\texttt{\$ sudo apt install zram-tools}
\end{highlightbox}

To check if it's installed:

\begin{highlightbox}
\texttt{\$ dpkg -l | grep zram}
\end{highlightbox}

To check if it's active:

\begin{highlightbox}
\texttt{\$ swapon --show}
\end{highlightbox}

\begin{infobox}
If you see something like \texttt{/dev/zram0}, it's working.
\end{infobox}

To check service status:

\begin{highlightbox}
\texttt{\$ systemctl status zramswap}
\end{highlightbox}

\section*{3 \diskon\ Understanding Your Disks}

Let's meet your disks.

\subsection*{\faTree\ \texttt{lsblk} — The Disk Tree Viewer}

\begin{highlightbox}
\texttt{\$ lsblk}
\end{highlightbox}

Example output:

\begin{highlightbox}
\texttt{sda}\\
\texttt{├─ sda1}\\
\texttt{├─ sda2}\\
\texttt{└─ sda3}\\
\texttt{sdb}\\
\texttt{└─ sdb1}
\end{highlightbox}

\subsubsection*{What it shows:}

\begin{itemize}
    \item Disk names (\texttt{sda}, \texttt{sdb})
    \item Partitions (\texttt{sda1}, \texttt{sda2})
    \item Mount points
    \item Size
    \item Type
\end{itemize}

More detailed output:

\begin{highlightbox}
\texttt{\$ lsblk -f}
\end{highlightbox}

This shows:

\begin{itemize}
    \item Filesystem type (ext4, vfat, ntfs)
    \item UUID
    \item Mountpoint
\end{itemize}

\begin{infobox}
\faFire\ Use this before and after plugging a USB.
\end{infobox}

\subsection*{\faKey\ \texttt{blkid} — UUID Finder}

\begin{highlightbox}
\texttt{\$ sudo blkid}
\end{highlightbox}

It shows something like:

\begin{highlightbox}
\texttt{/dev/sda1: UUID="1A2B-3C4D" TYPE="vfat"}
\end{highlightbox}

\begin{infobox}
\faQuestionCircle\ Why important?\\
Because in \texttt{/etc/fstab}, we use \textbf{UUID}, not \texttt{/dev/sda1}.\\
Because device names can change on boot. UUID never lies.
\end{infobox}

\section*{4 \mounticon\ Mounting and Unmounting Drives}

Unlike Windows, Linux does not auto-use disks unless mounted.

\subsection*{\faFolderOpen\ Manual Mount}

Step 1 — Create mount point:

\begin{highlightbox}
\texttt{\$ sudo mkdir /mnt/mydrive}
\end{highlightbox}

Step 2 — Mount:

\begin{highlightbox}
\texttt{\$ sudo mount /dev/sdb1 /mnt/mydrive}
\end{highlightbox}

Step 3 — Check:

\begin{highlightbox}
\texttt{\$ ls /mnt/mydrive}
\end{highlightbox}

\subsection*{\faEject\ Unmount}

\begin{highlightbox}
\texttt{\$ sudo umount /mnt/mydrive}
\end{highlightbox}

OR

\begin{highlightbox}
\texttt{\$ sudo umount /dev/sdb1}
\end{highlightbox}

\begin{warningbox}
\warningicon\ If busy error:
\begin{highlightbox}
\texttt{\$ lsof | grep sdb1}
\end{highlightbox}
\end{warningbox}

\section*{5 \faFile\ Making Mount Permanent — Editing \texttt{/etc/fstab}}

This file controls what mounts at boot.

Open it:

\begin{highlightbox}
\texttt{\$ sudo nano /etc/fstab}
\end{highlightbox}

Example entry:

\begin{highlightbox}
\texttt{UUID=1a2b3c4d-xxxx-xxxx-xxxx-xxxxxxxxxxxx  /mnt/mydrive  ext4  defaults  0  2}
\end{highlightbox}

\subsection*{Columns explained:}

\begin{infobox}
\begin{tabular}{ll}
\textbf{Column} & \textbf{Meaning} \\
UUID & Disk identifier \\
Mount point & Where it appears \\
Filesystem & ext4, ntfs, vfat \\
Options & defaults, noatime \\
Dump & usually 0 \\
fsck order & usually 2 \\
\end{tabular}
\end{infobox}

\subsection*{\faFlask\ Test Before Reboot}

After editing:

\begin{highlightbox}
\texttt{\$ sudo mount -a}
\end{highlightbox}

\begin{infobox}
\checkicon\ If no errors → safe to reboot.\\
\warningicon\ If error → FIX IT before rebooting.
\end{infobox}

\subsection*{Practical Example: Adding Second Drive}

\begin{infobox}
\arrowicon\ \href{https://mmbesar.github.io/Tutorials/2nd-drive/}{\texttt{mmbesar-2nd-drive}}
\end{infobox}

\section*{6 \ramicon\ Monitoring RAM Usage}

RAM = short-term memory.

\subsection*{\faBrain\ \texttt{free}}

\begin{highlightbox}
\texttt{\$ free -h}
\end{highlightbox}

Example:

\begin{highlightbox}
\texttt{              total   used   free  shared  buff/cache  available}\\
\texttt{Mem:           7.6G   3.1G   1.2G     500M     3.3G       3.8G}\\
\texttt{Swap:          2.0G   0B     2.0G}
\end{highlightbox}

Important columns:

\begin{itemize}
    \item \textbf{used} → currently used
    \item \textbf{available} → real usable memory
    \item \textbf{buff/cache} → Linux using RAM smartly
\end{itemize}

\begin{warningbox}
\warningicon\ Linux uses RAM aggressively. That's GOOD.\\
Unused RAM = wasted RAM.
\end{warningbox}

\subsection*{\faEye\ Real-time View}

\begin{highlightbox}
\texttt{\$ watch -n 1 free -h}
\end{highlightbox}

Updates every second.

\section*{7 \swapicon\ Understanding Swap}

Swap = disk used as overflow memory.

When RAM fills up → inactive pages move to swap.

But swap is slower than RAM.

Check swap:

\begin{highlightbox}
\texttt{\$ swapon --show}
\end{highlightbox}

or

\begin{highlightbox}
\texttt{\$ cat /proc/swaps}
\end{highlightbox}

\section*{8 \faFile\ Creating a Swap File (Manual Way)}

If no swap exists:

Step 1:

\begin{highlightbox}
\texttt{\$ sudo fallocate -l 2G /swapfile}
\end{highlightbox}

Step 2:

\begin{highlightbox}
\texttt{\$ sudo chmod 600 /swapfile}
\end{highlightbox}

Step 3:

\begin{highlightbox}
\texttt{\$ sudo mkswap /swapfile}
\end{highlightbox}

Step 4:

\begin{highlightbox}
\texttt{\$ sudo swapon /swapfile}
\end{highlightbox}

Make permanent:

Add to \texttt{/etc/fstab}:

\begin{highlightbox}
\texttt{/swapfile none swap sw 0 0}
\end{highlightbox}

\section*{9 \zramicon\ zram — Compressed RAM Swap}

zram creates compressed swap in RAM.

Instead of writing to disk:

\begin{itemize}
    \item It compresses pages
    \item Stores them in RAM
    \item Faster than disk swap
\end{itemize}

Great for:

\begin{itemize}
    \item 4GB RAM laptops
    \item Old machines
    \item Low-resource VPS
\end{itemize}

\subsection*{Install on Debian/Ubuntu/Mint}

\begin{highlightbox}
\texttt{\$ sudo apt install zram-tools}
\end{highlightbox}

Check config:

\begin{highlightbox}
\texttt{\$ cat /etc/default/zramswap}
\end{highlightbox}

You can adjust:

\begin{highlightbox}
\texttt{PERCENT=70}
\end{highlightbox}

\begin{infobox}
\faInfoCircle\ Meaning: use 70\% of RAM for zram.\\
You can change it but it is recommended to adjust it to [70\%-80\%].
\end{infobox}

Restart:

\begin{highlightbox}
\texttt{\$ sudo systemctl restart zramswap}
\end{highlightbox}

Check:

\begin{highlightbox}
\texttt{\$ swapon --show}
\end{highlightbox}

You should see \texttt{/dev/zram0}.

\begin{infobox}
\faStar\ Checkout This Practical Example in the link below:\\
\arrowicon\ \href{https://mmbesar.github.io/Tutorials/RAM-SWAP-ZRAM/}{\texttt{mmbesar-RAM-SWAP-ZRAM}}
\end{infobox}

\section*{10 \faBalanceScale\ Swap vs zram Comparison}

\begin{infobox}
\begin{tabular}{lll}
\textbf{Feature} & \textbf{Disk Swap} & \textbf{zram} \\
Speed & Slow & Faster \\
Uses disk & Yes & No \\
Best for & Large systems & Low RAM \\
Compression & No & Yes \\
\end{tabular}
\end{infobox}

\begin{infobox}
\faLightbulb\ Best practice for laptops:
\begin{itemize}
    \item Use zram
    \item Keep small disk swap
\end{itemize}
\end{infobox}

\section*{11 \faTools\ Troubleshooting Boot Failures (fstab Mistakes)}

If system won't boot:

\begin{enumerate}
    \item Boot into recovery mode
    \item Mount root
    \item Edit \texttt{/etc/fstab}
    \item Fix or comment broken line
\end{enumerate}

Pro tip:

Use this option in fstab for external drives:

\begin{highlightbox}
\texttt{nofail}
\end{highlightbox}

So system doesn't panic if drive missing.

\section*{12 \faList\ Useful Extra Commands}

Check disk usage:

\begin{highlightbox}
\texttt{\$ df -h}
\end{highlightbox}

Check disk space by folder:

\begin{highlightbox}
\texttt{\$ du -sh *}
\end{highlightbox}

See mount points:

\begin{highlightbox}
\texttt{\$ mount | grep sdb}
\end{highlightbox}

\section*{\labicon\ Practical Lab Exercises}

\begin{enumerate}
    \item \textbf{Exercise 1:} Plug USB → identify it using \texttt{lsblk}.
    \item \textbf{Exercise 2:} Mount it manually in \texttt{/mnt/testusb}.
    \item \textbf{Exercise 3:} Make it permanent in fstab.
    \item \textbf{Exercise 4:} Create 1GB swap file.
    \item \textbf{Exercise 5:} Install zram and compare swap usage.
\end{enumerate}

\section*{\videoicon\ Suggested Video References}

\begin{itemize}
    \item \href{https://www.youtube.com/watch?v=yWuHI7uoftY}{\texttt{https://www.youtube.com/watch?v=yWuHI7uoftY}}
    \item \href{https://youtu.be/CPvJvY79vEw}{\texttt{https://youtu.be/CPvJvY79vEw}}
\end{itemize}

\begin{infobox}
\faEye\ (You should watch after hands-on practice.)
\end{infobox}

\vspace{1em}
\begin{center}
\color{teamred}\rule{0.8\textwidth}{1pt}
\end{center}

\section*{\faFlagCheckered\ Final Advice}

Disk and memory management is not optional knowledge.

It's what separates:

\begin{center}
\textbf{"Linux user"}\\
\faArrowDown\\
\textbf{"Linux operator"} \rocketicon
\end{center}

\begin{itemize}
    \item Break things in a VM.
    \item Test everything with \texttt{mount -a}.
    \item Never reboot blindly after editing fstab.
\end{itemize}

\begin{infobox}
You're building serious Linux skills now.
\end{infobox}

% Updated Footer
\vspace{\fill}
\begin{center}
\color{muted}
All rights reserved for \texttt{TogetherTeam} [NOT For Commercial Use] !!\\
contact the author via email \href{mailto:khaledmhfz2004@gmail.com}{\texttt{khaledmhfz2004@gmail.com}}
\end{center}

\end{document}